% macros.tex — Vegas paper notation

% ── Packages ──────────────────────────────────────────────────────────
\usepackage{amsmath,amssymb,amsthm}
\usepackage{stmaryrd}       % \llbracket, \rrbracket
\usepackage{mathtools}       % \coloneqq
\usepackage{xcolor}
\usepackage{hyperref}
\usepackage{cleveref}
\usepackage{booktabs}
\usepackage{enumitem}
\usepackage{mathpartir}      % pretty inference rules (operational semantics)
\usepackage{simplebnf}       % pretty BNF grammars

% ── Theorem environments ──────────────────────────────────────────────
\newtheorem{theorem}{Theorem}[section]
\newtheorem{lemma}[theorem]{Lemma}
\newtheorem{proposition}[theorem]{Proposition}
\newtheorem{corollary}[theorem]{Corollary}
\theoremstyle{definition}
\newtheorem{definition}[theorem]{Definition}
\newtheorem{example}[theorem]{Example}
\theoremstyle{remark}
\newtheorem{remark}[theorem]{Remark}
\newtheorem{notation}[theorem]{Notation}

% ── Players and sets ──────────────────────────────────────────────────
\newcommand{\Players}{\mathcal{I}}

% ── Types and values ──────────────────────────────────────────────────
\newcommand{\Ty}{\mathsf{Ty}}
\newcommand{\Val}{\mathsf{Val}}

% ── Visibility annotations ────────────────────────────────────────────
\newcommand{\pub}{\ensuremath{\mathsf{public}}}
\newcommand{\sealed}{\ensuremath{\mathsf{sealed}}}

% ── Contexts ──────────────────────────────────────────────────────────
\newcommand{\view}[1]{\Gamma\!\!\downarrow_{#1}}       % viewVCtx
\newcommand{\erasectx}[1]{\| #1 \|}                    % eraseVCtx
\newcommand{\erasepub}[1]{\| #1 \|^{\pub}}             % erasePubVCtx (public-only erasure operator)

% ── Constructs ────────────────────────────────────────────────────────
\newcommand{\Ret}{\ensuremath{\mathbf{ret}}}
\newcommand{\Let}{\ensuremath{\mathbf{let}}}
\newcommand{\Sample}{\ensuremath{\mathbf{sample}}}
\newcommand{\Commit}{\ensuremath{\mathbf{commit}}}
\newcommand{\Reveal}{\ensuremath{\mathbf{reveal}}}
\newcommand{\Where}{\ensuremath{\mathbf{where}}}

% ── Small-step operational semantics ──────────────────────────────────
\newcommand{\sconf}[3]{\langle #1 \mathrel{;} #2 \mathrel{;} #3 \rangle} % configuration
\newcommand{\sstep}{\longrightarrow}                  % step relation
\newcommand{\lstep}[1]{\xrightarrow{\,#1\,}}          % labelled step
\newcommand{\sem}[1]{\llbracket #1 \rrbracket}        % evaluation in L
\newcommand{\supp}{\mathop{\mathrm{supp}}}            % support of a weighting
\newcommand{\outcomeof}{\mathsf{outcome}}             % terminal outcome reader

% ── Distributions ─────────────────────────────────────────────────────
\newcommand{\FWeight}[1]{\mathcal{W}(#1)}
\newcommand{\PMF}[1]{\mathrm{PMF}(#1)}

% ── Outcomes ──────────────────────────────────────────────────────────
\newcommand{\Outcome}{\mathcal{O}}

% ── Profiles and strategies ───────────────────────────────────────────
\newcommand{\Profile}{\sigma}

% ── Strategic bridge ──────────────────────────────────────────────────
\newcommand{\KG}{\ensuremath{\mathcal{G}}}
\newcommand{\EU}{\ensuremath{\mathrm{EU}}}

% ── Environments ──────────────────────────────────────────────────────
\newcommand{\Env}{\ensuremath{\mathrm{Env}}}

% ── Misc math ─────────────────────────────────────────────────────────
\newcommand{\NNQ}{\ensuremath{\mathbb{Q}_{\ge 0}}}
\newcommand{\bind}{\mathbin{\gg\!=}}

% ── Lean references ───────────────────────────────────────────────────
\newcommand{\TT}[1]{\texttt{#1}}

% ── Event graphs and configurations ───────────────────────────────────
\newcommand{\EG}{\mathcal{E}}                       % an event graph
\newcommand{\prereq}{\mathrm{prereq}}
\newcommand{\Cfg}{\mathsf{Cfg}}                     % configurations of an EG
\newcommand{\readsOf}{\mathrm{reads}}

% ── Machines ──────────────────────────────────────────────────────────
\newcommand{\State}{\mathsf{S}}
\newcommand{\Action}{\mathsf{A}}

% ── Solution concepts / kernels ───────────────────────────────────────
\newcommand{\Kbeh}{K^{\mathrm{beh}}}
\newcommand{\Kpure}{K^{\mathrm{pure}}}
\newcommand{\Kmix}{K^{\mathrm{mix}}}

% ── Faithfulness-relation shapes (the "forest" vocabulary) ────────────
% Symbols used in the guarantees map: an equivalence, a (one-directional)
% adequacy, and a refinement order.
\newcommand{\eqrel}{\cong}            % equivalence (symmetric)
\newcommand{\adeq}{\vDash}            % adequacy: right object IS the denotation
\newcommand{\ordrel}{\sqsubseteq}     % refinement order

% ── Non-intrusive "status / not-yet-mechanized" note ──────────────────
% A footnote that flags an aspirational or partial item without breaking
% the surrounding narrative.  Used sparingly.
\newcommand{\statusnote}[1]{\footnote{\textit{Status.}~#1}}

% ── Boxes for design notes / Lean correspondence ──────────────────────
\newcommand{\leancorr}[1]{%
  \par\smallskip
  {\raggedright\hbadness=10000\emergencystretch=3em
   \noindent\emph{Lean correspondence.} #1\par}%
  \smallskip}

% ── Lean signature boxes ──────────────────────────────────────────────
% A light-shaded, thin-ruled block for displaying Lean signatures.
% Contents are typeset (not verbatim) so we can mix \TT, math, and Greek
% as the rest of the document already does.
\usepackage[most]{tcolorbox}
\tcbset{
  leanstyle/.style={
    colback=black!3, colframe=black!30, boxrule=0.3pt, arc=1pt,
    left=8pt, right=8pt, top=4pt, bottom=4pt,
    breakable, enhanced jigsaw,
    before skip=4pt, after skip=4pt,
    fontupper=\footnotesize\ttfamily,
  },
}
\newtcolorbox{leanbox}[1][]{leanstyle, #1}
